\documentclass[
  doc,
  floatsintext,
  longtable,
  a4paper,
  nolmodern,
  notxfonts,
  notimes,
  12pt,
  colorlinks=true,linkcolor=blue,citecolor=blue,urlcolor=blue]{apa7}

\usepackage{amsmath}
\usepackage{amssymb}

\geometry{inner=1in, outer=1in}
\fancyhfoffset[LE,RO]{0cm}


\usepackage[bidi=default]{babel}
\babelprovide[main,import]{spanish}


% get rid of language-specific shorthands (see #6817):
\let\LanguageShortHands\languageshorthands
\def\languageshorthands#1{}

\RequirePackage{longtable}
\RequirePackage{threeparttablex}

\makeatletter
\renewcommand{\paragraph}{\@startsection{paragraph}{4}{\parindent}%
	{0\baselineskip \@plus 0.2ex \@minus 0.2ex}%
	{-.5em}%
	{\normalfont\normalsize\bfseries\typesectitle}}

\renewcommand{\subparagraph}[1]{\@startsection{subparagraph}{5}{0.5em}%
	{0\baselineskip \@plus 0.2ex \@minus 0.2ex}%
	{-\z@\relax}%
	{\normalfont\normalsize\bfseries\itshape\hspace{\parindent}{#1}\textit{\addperi}}{\relax}}
\makeatother




\usepackage{longtable, booktabs, multirow, multicol, colortbl, hhline, caption, array, float, xpatch}
\usepackage{subcaption}


\renewcommand\thesubfigure{\Alph{subfigure}}
\setcounter{topnumber}{2}
\setcounter{bottomnumber}{2}
\setcounter{totalnumber}{4}
\renewcommand{\topfraction}{0.85}
\renewcommand{\bottomfraction}{0.85}
\renewcommand{\textfraction}{0.15}
\renewcommand{\floatpagefraction}{0.7}

\usepackage{tcolorbox}
\tcbuselibrary{listings,theorems, breakable, skins}
\usepackage{fontawesome5}

\definecolor{quarto-callout-color}{HTML}{909090}
\definecolor{quarto-callout-note-color}{HTML}{0758E5}
\definecolor{quarto-callout-important-color}{HTML}{CC1914}
\definecolor{quarto-callout-warning-color}{HTML}{EB9113}
\definecolor{quarto-callout-tip-color}{HTML}{00A047}
\definecolor{quarto-callout-caution-color}{HTML}{FC5300}
\definecolor{quarto-callout-color-frame}{HTML}{ACACAC}
\definecolor{quarto-callout-note-color-frame}{HTML}{4582EC}
\definecolor{quarto-callout-important-color-frame}{HTML}{D9534F}
\definecolor{quarto-callout-warning-color-frame}{HTML}{F0AD4E}
\definecolor{quarto-callout-tip-color-frame}{HTML}{02B875}
\definecolor{quarto-callout-caution-color-frame}{HTML}{FD7E14}

%\newlength\Oldarrayrulewidth
%\newlength\Oldtabcolsep


\usepackage{hyperref}



\usepackage{color}
\usepackage{fancyvrb}
\newcommand{\VerbBar}{|}
\newcommand{\VERB}{\Verb[commandchars=\\\{\}]}
\DefineVerbatimEnvironment{Highlighting}{Verbatim}{commandchars=\\\{\}}
% Add ',fontsize=\small' for more characters per line
\usepackage{framed}
\definecolor{shadecolor}{RGB}{241,243,245}
\newenvironment{Shaded}{\begin{snugshade}}{\end{snugshade}}
\newcommand{\AlertTok}[1]{\textcolor[rgb]{0.68,0.00,0.00}{#1}}
\newcommand{\AnnotationTok}[1]{\textcolor[rgb]{0.37,0.37,0.37}{#1}}
\newcommand{\AttributeTok}[1]{\textcolor[rgb]{0.40,0.45,0.13}{#1}}
\newcommand{\BaseNTok}[1]{\textcolor[rgb]{0.68,0.00,0.00}{#1}}
\newcommand{\BuiltInTok}[1]{\textcolor[rgb]{0.00,0.23,0.31}{#1}}
\newcommand{\CharTok}[1]{\textcolor[rgb]{0.13,0.47,0.30}{#1}}
\newcommand{\CommentTok}[1]{\textcolor[rgb]{0.37,0.37,0.37}{#1}}
\newcommand{\CommentVarTok}[1]{\textcolor[rgb]{0.37,0.37,0.37}{\textit{#1}}}
\newcommand{\ConstantTok}[1]{\textcolor[rgb]{0.56,0.35,0.01}{#1}}
\newcommand{\ControlFlowTok}[1]{\textcolor[rgb]{0.00,0.23,0.31}{\textbf{#1}}}
\newcommand{\DataTypeTok}[1]{\textcolor[rgb]{0.68,0.00,0.00}{#1}}
\newcommand{\DecValTok}[1]{\textcolor[rgb]{0.68,0.00,0.00}{#1}}
\newcommand{\DocumentationTok}[1]{\textcolor[rgb]{0.37,0.37,0.37}{\textit{#1}}}
\newcommand{\ErrorTok}[1]{\textcolor[rgb]{0.68,0.00,0.00}{#1}}
\newcommand{\ExtensionTok}[1]{\textcolor[rgb]{0.00,0.23,0.31}{#1}}
\newcommand{\FloatTok}[1]{\textcolor[rgb]{0.68,0.00,0.00}{#1}}
\newcommand{\FunctionTok}[1]{\textcolor[rgb]{0.28,0.35,0.67}{#1}}
\newcommand{\ImportTok}[1]{\textcolor[rgb]{0.00,0.46,0.62}{#1}}
\newcommand{\InformationTok}[1]{\textcolor[rgb]{0.37,0.37,0.37}{#1}}
\newcommand{\KeywordTok}[1]{\textcolor[rgb]{0.00,0.23,0.31}{\textbf{#1}}}
\newcommand{\NormalTok}[1]{\textcolor[rgb]{0.00,0.23,0.31}{#1}}
\newcommand{\OperatorTok}[1]{\textcolor[rgb]{0.37,0.37,0.37}{#1}}
\newcommand{\OtherTok}[1]{\textcolor[rgb]{0.00,0.23,0.31}{#1}}
\newcommand{\PreprocessorTok}[1]{\textcolor[rgb]{0.68,0.00,0.00}{#1}}
\newcommand{\RegionMarkerTok}[1]{\textcolor[rgb]{0.00,0.23,0.31}{#1}}
\newcommand{\SpecialCharTok}[1]{\textcolor[rgb]{0.37,0.37,0.37}{#1}}
\newcommand{\SpecialStringTok}[1]{\textcolor[rgb]{0.13,0.47,0.30}{#1}}
\newcommand{\StringTok}[1]{\textcolor[rgb]{0.13,0.47,0.30}{#1}}
\newcommand{\VariableTok}[1]{\textcolor[rgb]{0.07,0.07,0.07}{#1}}
\newcommand{\VerbatimStringTok}[1]{\textcolor[rgb]{0.13,0.47,0.30}{#1}}
\newcommand{\WarningTok}[1]{\textcolor[rgb]{0.37,0.37,0.37}{\textit{#1}}}

\providecommand{\tightlist}{%
  \setlength{\itemsep}{0pt}\setlength{\parskip}{0pt}}
\usepackage{longtable,booktabs,array}
\newcounter{none} % for unnumbered tables
\usepackage{calc} % for calculating minipage widths
% Correct order of tables after \paragraph or \subparagraph
\usepackage{etoolbox}
\makeatletter
\patchcmd\longtable{\par}{\if@noskipsec\mbox{}\fi\par}{}{}
\makeatother
% Allow footnotes in longtable head/foot
\IfFileExists{footnotehyper.sty}{\usepackage{footnotehyper}}{\usepackage{footnote}}
\makesavenoteenv{longtable}

\usepackage{graphicx}
\makeatletter
\newsavebox\pandoc@box
\newcommand*\pandocbounded[1]{% scales image to fit in text height/width
  \sbox\pandoc@box{#1}%
  \Gscale@div\@tempa{\textheight}{\dimexpr\ht\pandoc@box+\dp\pandoc@box\relax}%
  \Gscale@div\@tempb{\linewidth}{\wd\pandoc@box}%
  \ifdim\@tempb\p@<\@tempa\p@\let\@tempa\@tempb\fi% select the smaller of both
  \ifdim\@tempa\p@<\p@\scalebox{\@tempa}{\usebox\pandoc@box}%
  \else\usebox{\pandoc@box}%
  \fi%
}
% Set default figure placement to htbp
\def\fps@figure{htbp}
\makeatother







\usepackage{newtx}

\defaultfontfeatures{Scale=MatchLowercase}
\defaultfontfeatures[\rmfamily]{Ligatures=TeX,Scale=1}





\title{Salida de texto con print()}


\shorttitle{SALIDA DE TEXTO CON PRINT()}


\usepackage{etoolbox}



\ccoppy{\textcopyright~2026}



\author{Edison Achalma}



\affiliation{
{Escuela Profesional de Economía, Universidad Nacional de San Cristóbal
de Huamanga}}




\leftheader{Achalma}

\date{2026-10-20}


\abstract{Este abstract será actualizado una vez que se complete el
contenido final del artículo. }



\authornote{\par{\addORCIDlink{Edison Achalma}{0000-0001-6996-3364}} 
\par{ }
\par{   El autor no tiene conflictos de interés que revelar.    Los
roles de autor se clasificaron utilizando la taxonomía de roles de
colaborador (CRediT; https://credit.niso.org/) de la siguiente
manera:  Edison Achalma:   conceptualización, metodología, análisis
formal, investigación, recursos, curación de
datos, redacción, visualización, supervisión, administración del
proyecto}
\par{La correspondencia relativa a este artículo debe dirigirse a Edison
Achalma, Escuela Profesional de Economía, Universidad Nacional de San
Cristóbal de Huamanga, Portal Independencia N°
57, Ayacucho, AYA 5001, Perú, Email: \href{mailto:elmer.achalma.09@unsch.edu.pe}{elmer.achalma.09@unsch.edu.pe}}
}

\makeatletter
\let\endoldlt\endlongtable
\def\endlongtable{
\hline
\endoldlt
}
\makeatother

\urlstyle{same}



\hypersetup{
  pdftitle={Salida de texto con print()},
  pdfauthor={Autor}
}
\makeatletter
\@ifpackageloaded{caption}{}{\usepackage{caption}}
\AtBeginDocument{%
\ifdefined\contentsname
  \renewcommand*\contentsname{Tabla de contenidos}
\else
  \newcommand\contentsname{Tabla de contenidos}
\fi
\ifdefined\listfigurename
  \renewcommand*\listfigurename{Lista de Figuras}
\else
  \newcommand\listfigurename{Lista de Figuras}
\fi
\ifdefined\listtablename
  \renewcommand*\listtablename{Lista de Tablas}
\else
  \newcommand\listtablename{Lista de Tablas}
\fi
\ifdefined\figurename
  \renewcommand*\figurename{Figura}
\else
  \newcommand\figurename{Figura}
\fi
\ifdefined\tablename
  \renewcommand*\tablename{Tabla}
\else
  \newcommand\tablename{Tabla}
\fi
}
\@ifpackageloaded{float}{}{\usepackage{float}}
\floatstyle{ruled}
\@ifundefined{c@chapter}{\newfloat{codelisting}{h}{lop}}{\newfloat{codelisting}{h}{lop}[chapter]}
\floatname{codelisting}{Listado}
\newcommand*\listoflistings{\listof{codelisting}{Lista de Listados}}
\makeatother
\makeatletter
\makeatother
\makeatletter
\@ifpackageloaded{caption}{}{\usepackage{caption}}
\@ifpackageloaded{subcaption}{}{\usepackage{subcaption}}
\makeatother
\makeatletter
\@ifpackageloaded{fontawesome5}{}{\usepackage{fontawesome5}}
\makeatother

% From https://tex.stackexchange.com/a/645996/211326
%%% apa7 doesn't want to add appendix section titles in the toc
%%% let's make it do it
\makeatletter
\xpatchcmd{\appendix}
  {\par}
  {\addcontentsline{toc}{section}{\@currentlabelname}\par}
  {}{}
\makeatother

%% Disable longtable counter
%% https://tex.stackexchange.com/a/248395/211326

\usepackage{etoolbox}

\makeatletter
\patchcmd{\LT@caption}
  {\bgroup}
  {\bgroup\global\LTpatch@captiontrue}
  {}{}
\patchcmd{\longtable}
  {\par}
  {\par\global\LTpatch@captionfalse}
  {}{}
\apptocmd{\endlongtable}
  {\ifLTpatch@caption\else\addtocounter{table}{-1}\fi}
  {}{}
\newif\ifLTpatch@caption
\makeatother

\begin{document}

\maketitle


\hypertarget{toc}{}
\tableofcontents
\newpage
\section[Introduction]{Salida de texto con print()}

\setcounter{secnumdepth}{3}

\setlength\LTleft{0pt}




\subsection{🎯 Motivación}\label{motivaciuxf3n}

\texttt{print()} es tu ventana al mundo: es cómo ves lo que tu programa
está haciendo. Sin \texttt{print()}, estarías programando a ciegas, sin
saber si tus variables tienen los valores correctos o si tu lógica
funciona. Dominar \texttt{print()} te permite ver, depurar y comunicar
resultados de manera efectiva.

\subsection{💡 Intuición}\label{intuiciuxf3n}

Imagina que \texttt{print()} es como un altavoz: le das un mensaje y lo
reproduce en voz alta para que todos lo escuchen. Puedes usarlo tantas
veces como quieras, y cada vez que lo usas, empieza una ``nueva línea''
automáticamente, como si presionaras Enter después de cada mensaje.

\subsection{🧠 Modelo mental}\label{modelo-mental}

La función \texttt{print()} funciona así:

\begin{enumerate}
\def\labelenumi{\arabic{enumi}.}
\tightlist
\item
  \textbf{Acepta uno o más argumentos}: puede recibir texto, números,
  variables, o combinaciones.
\item
  \textbf{Convierte todo a texto}: si le pasas un número, lo convierte a
  texto antes de mostrarlo.
\item
  \textbf{Añade un salto de línea al final}: por defecto, después de
  imprimir pasa a la siguiente línea.
\item
  \textbf{Puede personalizarse}: con parámetros como \texttt{end} y
  \texttt{sep} cambias su comportamiento.
\end{enumerate}

\subsection{📐 Desarrollo}\label{desarrollo}

\subsubsection{1. Sintaxis básica}\label{sintaxis-buxe1sica}

\begin{Shaded}
\begin{Highlighting}[]
\BuiltInTok{print}\NormalTok{(}\StringTok{"texto"}\NormalTok{)}
\BuiltInTok{print}\NormalTok{(variable)}
\BuiltInTok{print}\NormalTok{(}\StringTok{"texto"}\NormalTok{, variable)}
\end{Highlighting}
\end{Shaded}

\subsubsection{2. Comillas: simples o
dobles}\label{comillas-simples-o-dobles}

Python acepta ambos tipos de comillas. La única regla: \textbf{abre y
cierra con el mismo tipo}.

\begin{Shaded}
\begin{Highlighting}[]
\BuiltInTok{print}\NormalTok{(}\StringTok{"Hola"}\NormalTok{)      }\CommentTok{\# ✅ Válido}
\BuiltInTok{print}\NormalTok{(}\StringTok{\textquotesingle{}Hola\textquotesingle{}}\NormalTok{)      }\CommentTok{\# ✅ Válido}
\BuiltInTok{print}\NormalTok{(}\StringTok{"Hola\textquotesingle{})      \# ❌ Error: mezcla de comillas}
\end{Highlighting}
\end{Shaded}

\textbf{¿Cuándo usar cada una?}

\begin{itemize}
\tightlist
\item
  Si tu texto contiene comillas simples, usa dobles:
  \texttt{print("It\textquotesingle{}s\ awesome")}
\item
  Si contiene comillas dobles, usa simples:
  \texttt{print(\textquotesingle{}Ella\ dijo\ "Hola"\textquotesingle{})}
\item
  O usa \texttt{\textbackslash{}} para escapar:
  \texttt{print("It\textbackslash{}\textquotesingle{}s\ awesome")}
\end{itemize}

\subsubsection{\texorpdfstring{3. Múltiples \texttt{print()}: cada uno
en nueva
línea}{3. Múltiples print(): cada uno en nueva línea}}\label{muxfaltiples-print-cada-uno-en-nueva-luxednea}

Por defecto, cada \texttt{print()} comienza en una línea nueva:

\begin{Shaded}
\begin{Highlighting}[]
\BuiltInTok{print}\NormalTok{(}\StringTok{"Línea 1"}\NormalTok{)}
\BuiltInTok{print}\NormalTok{(}\StringTok{"Línea 2"}\NormalTok{)}
\BuiltInTok{print}\NormalTok{(}\StringTok{"Línea 3"}\NormalTok{)}
\CommentTok{\# Salida:}
\CommentTok{\# Línea 1}
\CommentTok{\# Línea 2}
\CommentTok{\# Línea 3}
\end{Highlighting}
\end{Shaded}

\subsubsection{\texorpdfstring{4. Múltiples argumentos en un
\texttt{print()}}{4. Múltiples argumentos en un print()}}\label{muxfaltiples-argumentos-en-un-print}

Puedes pasar varios elementos separados por comas:

\begin{Shaded}
\begin{Highlighting}[]
\NormalTok{nombre }\OperatorTok{=} \StringTok{"Ana"}
\NormalTok{edad }\OperatorTok{=} \DecValTok{25}
\BuiltInTok{print}\NormalTok{(}\StringTok{"Nombre:"}\NormalTok{, nombre, }\StringTok{"Edad:"}\NormalTok{, edad)}
\CommentTok{\# Salida: Nombre: Ana Edad: 25}
\end{Highlighting}
\end{Shaded}

\subsubsection{\texorpdfstring{5. Personalizar el final:
\texttt{end}}{5. Personalizar el final: end}}\label{personalizar-el-final-end}

El parámetro \texttt{end} controla qué se añade al final (por defecto es
\texttt{\textbackslash{}n}, salto de línea):

\begin{Shaded}
\begin{Highlighting}[]
\BuiltInTok{print}\NormalTok{(}\StringTok{"Hola"}\NormalTok{, end}\OperatorTok{=}\StringTok{" "}\NormalTok{)}
\BuiltInTok{print}\NormalTok{(}\StringTok{"Mundo"}\NormalTok{)}
\CommentTok{\# Salida: Hola Mundo}

\BuiltInTok{print}\NormalTok{(}\StringTok{"Manzana"}\NormalTok{, end}\OperatorTok{=}\StringTok{", "}\NormalTok{)}
\BuiltInTok{print}\NormalTok{(}\StringTok{"Pera"}\NormalTok{, end}\OperatorTok{=}\StringTok{", "}\NormalTok{)}
\BuiltInTok{print}\NormalTok{(}\StringTok{"Banana"}\NormalTok{)}
\CommentTok{\# Salida: Manzana, Pera, Banana}
\end{Highlighting}
\end{Shaded}

\subsubsection{\texorpdfstring{6. Personalizar el separador:
\texttt{sep}}{6. Personalizar el separador: sep}}\label{personalizar-el-separador-sep}

El parámetro \texttt{sep} controla qué se pone entre múltiples
argumentos (por defecto es espacio):

\begin{Shaded}
\begin{Highlighting}[]
\BuiltInTok{print}\NormalTok{(}\StringTok{"uno"}\NormalTok{, }\StringTok{"dos"}\NormalTok{, }\StringTok{"tres"}\NormalTok{, sep}\OperatorTok{=}\StringTok{"{-}"}\NormalTok{)}
\CommentTok{\# Salida: uno{-}dos{-}tres}

\BuiltInTok{print}\NormalTok{(}\DecValTok{1}\NormalTok{, }\DecValTok{2}\NormalTok{, }\DecValTok{3}\NormalTok{, sep}\OperatorTok{=}\StringTok{" + "}\NormalTok{)}
\CommentTok{\# Salida: 1 + 2 + 3}
\end{Highlighting}
\end{Shaded}

\subsection{🔬 Ejemplos}\label{ejemplos}

\textbf{Mínimo ---} Texto simple:

\begin{Shaded}
\begin{Highlighting}[]
\BuiltInTok{print}\NormalTok{(}\StringTok{"Hola Mundo"}\NormalTok{)}
\CommentTok{\# Salida: Hola Mundo}
\end{Highlighting}
\end{Shaded}

\textbf{Progresivo ---} Variaciones de comillas:

\begin{Shaded}
\begin{Highlighting}[]
\BuiltInTok{print}\NormalTok{(}\StringTok{\textquotesingle{}Python es "genial"\textquotesingle{}}\NormalTok{)}
\BuiltInTok{print}\NormalTok{(}\StringTok{"It\textquotesingle{}s a beautiful day"}\NormalTok{)}
\BuiltInTok{print}\NormalTok{(}\StringTok{"Ella dijo: }\CharTok{\textbackslash{}"}\StringTok{Hola}\CharTok{\textbackslash{}"}\StringTok{"}\NormalTok{)  }\CommentTok{\# Escapando comillas}
\CommentTok{\# Salida:}
\CommentTok{\# Python es "genial"}
\CommentTok{\# It\textquotesingle{}s a beautiful day}
\CommentTok{\# Ella dijo: "Hola"}
\end{Highlighting}
\end{Shaded}

\textbf{Real ---} Menú de opciones con formato:

\begin{Shaded}
\begin{Highlighting}[]
\BuiltInTok{print}\NormalTok{(}\StringTok{"=== MENÚ PRINCIPAL ==="}\NormalTok{)}
\BuiltInTok{print}\NormalTok{(}\StringTok{"1. Ver perfil"}\NormalTok{)}
\BuiltInTok{print}\NormalTok{(}\StringTok{"2. Editar configuración"}\NormalTok{, end}\OperatorTok{=}\StringTok{" "}\NormalTok{)}
\BuiltInTok{print}\NormalTok{(}\StringTok{"(requiere permisos)"}\NormalTok{)}
\BuiltInTok{print}\NormalTok{(}\StringTok{"3. Salir"}\NormalTok{)}
\BuiltInTok{print}\NormalTok{(}\StringTok{"======================"}\NormalTok{)}
\CommentTok{\# Salida:}
\CommentTok{\# === MENÚ PRINCIPAL ===}
\CommentTok{\# 1. Ver perfil}
\CommentTok{\# 2. Editar configuración (requiere permisos)}
\CommentTok{\# 3. Salir}
\CommentTok{\# ======================}
\end{Highlighting}
\end{Shaded}

\subsection{⚠️ Errores comunes}\label{errores-comunes}

{\def\LTcaptype{none} % do not increment counter
\begin{longtable}[]{@{}
  >{\raggedright\arraybackslash}p{(\linewidth - 4\tabcolsep) * \real{0.2292}}
  >{\raggedright\arraybackslash}p{(\linewidth - 4\tabcolsep) * \real{0.3681}}
  >{\raggedright\arraybackslash}p{(\linewidth - 4\tabcolsep) * \real{0.4028}}@{}}
\toprule\noalign{}
\begin{minipage}[b]{\linewidth}\raggedright
Error
\end{minipage} & \begin{minipage}[b]{\linewidth}\raggedright
Síntoma
\end{minipage} & \begin{minipage}[b]{\linewidth}\raggedright
Cómo detectarlo
\end{minipage} \\
\midrule\noalign{}
\endhead
\bottomrule\noalign{}
\endlastfoot
\textbf{Olvidar comillas} & \texttt{SyntaxError:\ invalid\ syntax} & El
texto siempre va entre \texttt{"\ "} o
\texttt{\textquotesingle{}\ \textquotesingle{}} \\
\textbf{Comillas desbalanceadas} &
\texttt{SyntaxError:\ unterminated\ string\ literal} & Cuenta que abres
y cierras el mismo número de comillas \\
\textbf{Confundir \texttt{print} con \texttt{Print}} &
\texttt{NameError:\ name\ \textquotesingle{}Print\textquotesingle{}\ is\ not\ defined}
& Python es sensible a mayúsculas: usa \texttt{print} en minúsculas \\
\textbf{Olvidar paréntesis} &
\texttt{SyntaxError:\ Missing\ parentheses\ in\ call\ to\ \textquotesingle{}print\textquotesingle{}}
& Python 3 requiere \texttt{()} \\
\textbf{Espacio extra al final} & No da error, pero puede afectar
formato & Usa \texttt{end=""} sin espacio si no quieres separador \\
\end{longtable}
}

\subsection{✅ Buenas prácticas}\label{buenas-pruxe1cticas}

\begin{itemize}
\item
  \textbf{Usa f-strings} para combinar texto y variables (lo veremos
  pronto):

\begin{Shaded}
\begin{Highlighting}[]
\NormalTok{nombre }\OperatorTok{=} \StringTok{"Ana"}
\BuiltInTok{print}\NormalTok{(}\SpecialStringTok{f"Hola, }\SpecialCharTok{\{}\NormalTok{nombre}\SpecialCharTok{\}}\SpecialStringTok{"}\NormalTok{)  }\CommentTok{\# Más legible que concatenar}
\end{Highlighting}
\end{Shaded}
\item
  \textbf{Comillas simples por defecto}, usa dobles solo cuando el texto
  contenga comillas simples.
\item
  \textbf{\texttt{end="\ "}} con espacio para mejor legibilidad al unir
  líneas.
\item
  \textbf{No concatenes con \texttt{+}} innecesariamente: usa múltiples
  argumentos en \texttt{print()} con comas.
\item
  \textbf{Mensajes descriptivos}: en lugar de \texttt{print(x)}, usa
  \texttt{print("El\ resultado\ es:",\ x)}.
\end{itemize}

\subsection{📝 Ejercicio}\label{ejercicio}

\textbf{Básico} --- Formato de salida

Escribe un programa que:

\begin{enumerate}
\def\labelenumi{\arabic{enumi}.}
\tightlist
\item
  Imprima ``--- BIENVENIDO ---'' en una línea
\item
  Imprima ``Usuario:'' y luego el nombre ``Carlos'' en la misma línea,
  separados por un espacio
\item
  Imprima ``Edad:'' y luego ``30'' en la misma línea
\item
  Imprima ``Estado: Activo'' con un guión bajo (\texttt{\_}) entre las
  palabras
\item
  Imprima ``--- FIN ---'' con un punto final (\texttt{.}) al final, sin
  salto de línea
\end{enumerate}

\textbf{Pista:} usa los parámetros \texttt{sep}, \texttt{end} y
múltiples argumentos.

\begin{center}\rule{0.5\linewidth}{0.5pt}\end{center}

\textbf{Solución (no mires hasta haberlo intentado):}

\begin{Shaded}
\begin{Highlighting}[]
\BuiltInTok{print}\NormalTok{(}\StringTok{"{-}{-}{-} BIENVENIDO {-}{-}{-}"}\NormalTok{)}
\BuiltInTok{print}\NormalTok{(}\StringTok{"Usuario:"}\NormalTok{, }\StringTok{"Carlos"}\NormalTok{, sep}\OperatorTok{=}\StringTok{" "}\NormalTok{)}
\BuiltInTok{print}\NormalTok{(}\StringTok{"Edad:"}\NormalTok{, }\DecValTok{30}\NormalTok{, sep}\OperatorTok{=}\StringTok{" "}\NormalTok{)}
\BuiltInTok{print}\NormalTok{(}\StringTok{"Estado"}\NormalTok{, }\StringTok{"Activo"}\NormalTok{, sep}\OperatorTok{=}\StringTok{"\_"}\NormalTok{)}
\BuiltInTok{print}\NormalTok{(}\StringTok{"{-}{-}{-} FIN {-}{-}{-}"}\NormalTok{, end}\OperatorTok{=}\StringTok{"."}\NormalTok{)}
\end{Highlighting}
\end{Shaded}

\textbf{Salida esperada:}

\begin{verbatim}
--- BIENVENIDO ---
Usuario: Carlos
Edad: 30
Estado_Activo
--- FIN ---.
\end{verbatim}

\subsection{🗂️ Síntesis}\label{suxedntesis}

\begin{itemize}
\tightlist
\item
  \textbf{\texttt{print()}} muestra texto y valores en la consola.
\item
  \textbf{Comillas simples o dobles} son válidas, pero deben coincidir.
\item
  \textbf{Cada \texttt{print()}} añade un salto de línea por defecto.
\item
  \textbf{\texttt{end}} personaliza el final de línea (por defecto
  \texttt{\textbackslash{}n}).
\item
  \textbf{\texttt{sep}} personaliza el separador entre argumentos (por
  defecto espacio).
\item
  \textbf{Múltiples argumentos} separados por comas se imprimen con
  espacio entre ellos.
\end{itemize}

\subsection{🔗 Qué sigue}\label{quuxe9-sigue}

\begin{itemize}
\tightlist
\item
  \textbf{Relacionado:} variables (cómo almacenar datos para imprimir
  después), f-strings (formato avanzado)
\item
  \textbf{Siguiente tema:} Variables y tipos de datos --- aprenderás a
  crear y manejar datos que luego podrás imprimir.
\end{itemize}

\section{Publicaciones Similares}\label{publicaciones-similares}

Si te interesó este artículo, te recomendamos que explores otros blogs y
recursos relacionados que pueden ampliar tus conocimientos. Aquí te dejo
algunas sugerencias:

\begin{enumerate}
\def\labelenumi{\arabic{enumi}.}
\tightlist
\item
  \href{https://numerus-scriptum.netlify.app/python/2020-06-19-instalacion-de-anaconda/index.pdf}{\faIcon{file-pdf}}
  \href{https://numerus-scriptum.netlify.app/python/2020-06-19-instalacion-de-anaconda}{Instalacion
  De Anaconda}
\item
  \href{https://numerus-scriptum.netlify.app/python/2020-06-20-configurar-entorno-virtual-python-anaconda/index.pdf}{\faIcon{file-pdf}}
  \href{https://numerus-scriptum.netlify.app/python/2020-06-20-configurar-entorno-virtual-python-anaconda}{Configurar
  Entorno Virtual Python Anaconda}
\item
  \href{https://numerus-scriptum.netlify.app/python/2021-04-17-01-introducion-a-la-programacion-con-python/index.pdf}{\faIcon{file-pdf}}
  \href{https://numerus-scriptum.netlify.app/python/2021-04-17-01-introducion-a-la-programacion-con-python}{01
  Introducion A La Programacion Con Python}
\item
  \href{https://numerus-scriptum.netlify.app/python/2021-05-31-02-variables-expresiones-y-statements-con-python/index.pdf}{\faIcon{file-pdf}}
  \href{https://numerus-scriptum.netlify.app/python/2021-05-31-02-variables-expresiones-y-statements-con-python}{02
  Variables Expresiones Y Statements Con Python}
\item
  \href{https://numerus-scriptum.netlify.app/python/2021-06-07-03-objetos-de-python/index.pdf}{\faIcon{file-pdf}}
  \href{https://numerus-scriptum.netlify.app/python/2021-06-07-03-objetos-de-python}{03
  Objetos De Python}
\item
  \href{https://numerus-scriptum.netlify.app/python/2021-06-14-04-ejecucion-condicional-con-python/index.pdf}{\faIcon{file-pdf}}
  \href{https://numerus-scriptum.netlify.app/python/2021-06-14-04-ejecucion-condicional-con-python}{04
  Ejecucion Condicional Con Python}
\item
  \href{https://numerus-scriptum.netlify.app/python/2021-06-21-05-iteraciones-con-python/index.pdf}{\faIcon{file-pdf}}
  \href{https://numerus-scriptum.netlify.app/python/2021-06-21-05-iteraciones-con-python}{05
  Iteraciones Con Python}
\item
  \href{https://numerus-scriptum.netlify.app/python/2021-08-16-06-funciones-con-python/index.pdf}{\faIcon{file-pdf}}
  \href{https://numerus-scriptum.netlify.app/python/2021-08-16-06-funciones-con-python}{06
  Funciones Con Python}
\item
  \href{https://numerus-scriptum.netlify.app/python/2021-08-23-07-dataframes-con-python/index.pdf}{\faIcon{file-pdf}}
  \href{https://numerus-scriptum.netlify.app/python/2021-08-23-07-dataframes-con-python}{07
  Dataframes Con Python}
\item
  \href{https://numerus-scriptum.netlify.app/python/2021-11-29-08-prediccion-y-metrica-de-performance-con-python/index.pdf}{\faIcon{file-pdf}}
  \href{https://numerus-scriptum.netlify.app/python/2021-11-29-08-prediccion-y-metrica-de-performance-con-python}{08
  Prediccion Y Metrica De Performance Con Python}
\item
  \href{https://numerus-scriptum.netlify.app/python/2021-12-06-09-metodos-de-machine-learning-para-clasificacion-con-python/index.pdf}{\faIcon{file-pdf}}
  \href{https://numerus-scriptum.netlify.app/python/2021-12-06-09-metodos-de-machine-learning-para-clasificacion-con-python}{09
  Metodos De Machine Learning Para Clasificacion Con Python}
\item
  \href{https://numerus-scriptum.netlify.app/python/2021-12-13-10-metodos-de-machine-learning-para-regresion-con-python/index.pdf}{\faIcon{file-pdf}}
  \href{https://numerus-scriptum.netlify.app/python/2021-12-13-10-metodos-de-machine-learning-para-regresion-con-python}{10
  Metodos De Machine Learning Para Regresion Con Python}
\item
  \href{https://numerus-scriptum.netlify.app/python/2025-05-10-visualizacion-de-datos-con-python/index.pdf}{\faIcon{file-pdf}}
  \href{https://numerus-scriptum.netlify.app/python/2025-05-10-visualizacion-de-datos-con-python}{Visualizacion
  De Datos Con Python}
\end{enumerate}

Esperamos que encuentres estas publicaciones igualmente interesantes y
útiles. ¡Disfruta de la lectura!






\end{document}
